%% kapa: exam=1 round=36 year=2025 period=2 form=A subject=accounting subject_ko=회계학 seq=33 number=73
%% source: Q-Net 감정평가사 기출문제 (cst00309) artlSeq=5231525 file=raw/1008/5231525/2025년 제36회 감정평가사 1차시험 2교시 문제지 원본.pdf
%% conversion: kordoc 4.24.51 -> tools/gen1.py (auto); review=proofread-sonnet; flags=none
\begin{question}{73}
(주)감평은 단일제품을 대량생산하고 있으며, 선입선출법에 의한 종합원가계산을 적용하고 있다. 직접재료원가는 공정 초에 전량 투입되며, 전환원가(conversion costs, 가공원가)는 공정 전반에 걸쳐 균등하게 발생한다. 품질검사는 생산공정의 50 \% 시점에서 이루어지며, 당기 품질검사를 통과한 합격품의 3 \%를 정상공손으로 간주한다.
(주)감평의 20×1년의 생산 및 원가와 관련한 자료는 다음과 같다.
\begin{tcolorbox}[enhanced, colback=white, colframe=black, boxrule=0.4pt, arc=0pt, left=4pt, right=4pt, top=4pt, bottom=4pt]
\begin{tabular}{@{}l r c r r@{}}
 & \multicolumn{1}{c}{\uline{\hspace{1em}수량\hspace{1em}}} & \uline{\hspace{1em}완성도\hspace{1em}} & \multicolumn{1}{c}{\uline{\hspace{1em}직접재료원가\hspace{1em}}} & \multicolumn{1}{c}{\uline{\hspace{1em}전환원가\hspace{1em}}} \\
기초재공품 & 500단위 & 40 \% & ￦32,000 & ￦18,380 \\
당기착수량 & 8,500단위 & & 850,000 & 409,000 \\
당기완성량 & 8,000단위 & & & \\
기말재공품 & 600단위 & 30 \% & &
\end{tabular}
\end{tcolorbox}
(주)감평의 20×1년 정상공손원가는?
\begin{choices}
  \item ￦28,125
  \item ￦30,000
  \item ￦30,375
  \item ￦32,250
  \item ￦35,000
\end{choices}
\end{question}
\answer{2}
